\documentclass[10pt,a4paper]{amsart}
\usepackage[margin=19mm]{geometry}
\usepackage{amsmath,amssymb,mathtools}
\usepackage[scr=rsfs,cal=euler]{mathalfa}
\usepackage{xfrac}
\usepackage[unicode,hidelinks,bookmarks=false]{hyperref}
\newcommand{\ie}{i.e.\ }
\newcommand{\cf}{cf.\ }
\newcommand{\resp}{respectively\ }
\newcommand{\Subsection}{\subsection}
\providecommand{\nobreakdash}{\nobreak\hbox{-}}
\setlength{\emergencystretch}{2em}
\multlinegap0pt
\usepackage[shortlabels]{enumitem}
\usepackage{amssymb}
\usepackage{xcolor}
\usepackage{algorithm}\usepackage{algpseudocode}
\usepackage{tikz}
\newtheorem{lemma}{Lemma}
\newcommand\C{\mathcal{C}}
\newcommand\argmin{\mathop{\rm argmin}}
\title{Convergence of the Frank-Wolfe Algorithm for Monotone Variational Inequalities}
\author{Matthew Hough}
\date{Source: arXiv:2603.12104v2 (CC BY 4.0)}
\begin{document}
\maketitle

\noindent\emph{Source and licence.} Convergence of the Frank-Wolfe Algorithm for Monotone Variational Inequalities. Version arXiv:2603.12104v2 (CC BY 4.0), distributed by arXiv under the Creative Commons Attribution 4.0 International licence, http://creativecommons.org/licenses/by/4.0/, as recorded in the arXiv licence metadata for this exact version and shown on the abstract page https://arxiv.org/abs/2603.12104v2. Full licence notice: \texttt{licenses/arxiv-2603.12104-CC-BY-4.0.txt}.\par\smallskip

\noindent\textit{Editorial scope.} Counted unit: the article's lemma lem:lmo-lipschitz (source0.tex lines 696--709). The article's complete original proof of it is delivered with it (source0.tex lines 710--761, one \texttt{proof} environment of 52 lines). No mathematics is written, repaired or re-derived: the statement and the proof are the article's own lines, in the article's own notation, and no retained line is substituted or re-flowed.  No further source-local context is needed: every cross-reference of every retained line resolves inside the two delivered spans.  Nothing was omitted from any retained span, so the package carries no omission record.  No retained line contains a citation, so the wrapper carries no bibliography and no bibliography entry.  No retained line contains a figure environment, an image-inclusion command or drawing code, so no image is delivered and the package declares no asset dependency.  Editorial formatting only, in addition: an AMS \texttt{amsart} preamble that restates the article's own declarations and macros used by the retained lines, namely the environments \texttt{lemma} on separate counters, together with the standard packages those lines need.  The retained environments print in this document as Lemma 1 in order of appearance, whereas the article prints the same items with the numbering of its own full document.  The complete native source of the article as distributed in the arXiv e-print for this exact version is bundled unchanged as \texttt{sources/196260/source0.tex}.
\begin{lemma} \label{lem:lmo-lipschitz}
  Let $\C \subseteq \mathbb{R}^n$ be nonempty, compact, and $\alpha$-strongly convex for some $\alpha > 0$. Recall the definition of the linear minimization oracle
  \[
    \beta(u) := \argmin_{s\in\C}\langle{u, s}\rangle,\quad u \in \mathbb{R}^n.
  \]
  Then,
  \begin{enumerate}[label=(\roman*)]
      \item For every $u \in \mathbb{R}^n\setminus\{0\}$, the minimizer $\beta(u)$ is unique.
      \item The function $\beta(u)$ is $(1/\alpha)$-Lipschitz on the unit sphere; that is, for all unit vectors $u,v \in \mathbb{R}^n$
      \begin{equation} \label{eq:lmo-lip}
          \lVert \beta(u) - \beta(v)\rVert \leq \frac{1}{\alpha}\lVert u - v\rVert.
      \end{equation}
  \end{enumerate}
\end{lemma}

\begin{proof}
  First, since $\C$ is nonempty and compact, $\beta(u)$ must exist. Now suppose $u \neq 0$ and there exist distinct $s_1,s_2 \in \C$ with
  $\langle{u,s_1}\rangle = \langle{u,s_2}\rangle = k \in \mathbb{R}$. Define $m = (s_1 + s_2)/2$. By $\alpha$-strong convexity of $\C$, for $t=1/2$ we have
  \[
    B(m,\rho) \subseteq \C,
  \]
  where $\rho = \frac{\alpha}{8}\lVert s_1 - s_2\rVert^2$, which is positive by the assumption that $s_1 \neq s_2$. Take $w = u/\lVert u\rVert$ and
  $z = m - \rho w$. We must have $z \in \C$, while
  \begin{align*}
    \langle{u,z}\rangle &= \langle{u,m}\rangle - \rho\langle{u,w}\rangle\\
                        &= k - \rho\lVert u\rVert\\
                        &< k.
  \end{align*}
  But we have now found a $z\in \C$ which contradicts the minimality of the claimed minimizers. It can only be that $\rho = 0$, which implies $s_1 = s_2$, completing the proof of (i).

  To prove (ii), fix unit vectors $w_1, w_2 \in \mathbb{R}^n$ and let $s_1 := \beta(w_1), s_2 := \beta(w_2)$, $d := \lVert s_1 - s_2\rVert$.
  If $d = 0$ the proof is trivial, so suppose $d > 0$.
  For any $t \in [0,1]$, strong convexity of $\C$ implies $B(m_t, r_t) \subseteq \C$, where
  $m_t = (1-t)s_1 + ts_2$ and $r_t = \frac{\alpha}{2}t(1-t)d^2$. Take $z_t := m_t - r_t w_1$, which must lie in $\C$.
  By optimality of $s_1$ for $\min_{s\in \C}\langle{w_1,s}\rangle$,
  \[
    \langle{w_1,s_1}\rangle \leq \langle{w_1, z_t}\rangle = (1-t)\langle{w_1,s_1}\rangle + t\langle{w_1,s_2}\rangle - r_t,
  \]
  which after rearranging gives
  \[
    t\langle{w_1, s_2 - s_1}\rangle \geq r_t = \frac{\alpha}{2}t(1-t)d^2.
  \]
  So for all $t \in (0,1)$,
  \[
    \langle{w_1,s_2 - s_1}\rangle \geq \frac{\alpha}{2}(1-t)d^2.
  \]
  Taking the limit as $t \downarrow 0$, 
  \begin{equation} \label{eq:int-support-l-bound-1}
    \langle{w_1,s_2 - s_1}\rangle \geq \frac{\alpha}{2}d^2.
  \end{equation}
  Swapping $w_1$ and $w_2$ and $s_1$ and $s_2$ in the above working gives the symmetric form
  \begin{equation} \label{eq:int-support-l-bound-2}
    \langle{w_2,s_1 - s_2}\rangle \geq \frac{\alpha}{2}d^2.
  \end{equation}
  Adding \eqref{eq:int-support-l-bound-1} and \eqref{eq:int-support-l-bound-2} we get
  \[
    \langle{w_1 - w_2, s_2 - s_1}\rangle \geq \alpha d^2,
  \]
  and by Cauchy-Schwarz
  \[
    \alpha d^2 \leq \lVert w_1 - w_2\rVert d,
  \]
  which gives the desired inequality
  \[
    \lVert \beta(w_1) - \beta(w_2)\rVert \leq \frac{1}{\alpha}\lVert w_1 - w_2\rVert.
  \]
\end{proof}

\end{document}
